\section{Conclusion}
\label{sec:conclusion-ra}

The amendment side of constitutional admission is type-conditioned on
a backward-refinement witness; the response side, for actions whose
effects are reversible by an inverse executor, can be conditioned on
a positive TTL witness and an optional rollback receipt that closes
the action chain. The composition theorem
\thm{ttl_bounded_amendment_chain_preserves_baseline} relates the
response side to the amendment side: a chain of TTL-bounded
amendments, each carrying a per-step refinement witness, preserves
the baseline constitution at every instant under the auto-reversion
semantics of the TTL window. The reduction
\thm{destructive_action_requires_bilateral_admission} specializes the
parent paper's treaty intersection theorem to the destructive
subclass of the action taxonomy, making human-in-the-loop oversight
of destructive enforcement a typed obligation rather than a policy
file. The endpoint deployment instance grounds the substrate in four
reversible action variants with content-hash-gated inverse executors
and a four-receipt grammar; the TTL auto-expiry scheduler, the
missing inverse executors on the destructive subclass, and the
bilateral cosignature path remain substrate-level obligations the
present construction does not deploy. A focused Lean session on the
headline composition theorem is the next falsifiable step: if the
theorem discharges to \emph{rfl}, the construction collapses to the
definitional bridges; if it discharges by induction over the chain
with per-step witnesses, the response side reaches the same type-
level discipline as the amendment side of the parent paper.
